% 奥利弗·亥维赛（综述）
% license CCBYSA3
% type Wiki

本文根据 CC-BY-SA 协议转载翻译自维基百科\href{https://en.wikipedia.org/wiki/Oliver_Heaviside}{相关文章}

奥利弗·赫维赛德（Oliver Heaviside，/ˈhɛvisaɪd/，读作“赫维赛德”\(^\text{[2]}\)，1850年5月18日－1925年2月3日）是一位英国数学家和电气工程师。他发明了一种求解微分方程的新方法（与拉普拉斯变换等价），独立发展了向量分析，并将麦克斯韦方程组改写成如今通用的形式。他对麦克斯韦方程组在麦克斯韦逝世后的理解与应用产生了深远影响。此外，1893年，赫维赛德将麦克斯韦方程扩展到了“引力电磁学”领域——这一理论在2005年的“重力探测器B”实验中得到了验证。他提出的“电报方程”在他生前的晚期获得了商业上的重要意义，尽管这一成果曾长期未被重视，因为当时很少有人理解他所采用的创新数学方法。\(^\text{[3]}\)尽管赫维赛德一生大部分时间都与主流科学界格格不入，但他彻底改变了电信、数学和科学的面貌。\(^\text{[3]}\)
\subsection{早年经历}
奥利弗·赫维赛德于1850年5月18日出生在英国伦敦卡姆登镇国王街55号（现为普伦德街，Plender Street）\(^\text{[4]: 13 }\)，是托马斯·赫维赛德和雷切尔·伊丽莎白·韦斯特的第三个孩子，也是家中最小的儿子。他的父亲是一名制图员和木版雕刻师。赫维赛德从小个子矮小、头发红亮，年幼时患上猩红热，导致听力受损。家里继承了一小笔遗产，使得全家在他十三岁时搬到了卡姆登较好的地区，他也因此得以进入卡姆登文法学校就读。赫维赛德是个出色的学生，1865年在500名学生中名列第五。然而，由于家庭经济拮据，他在16岁时被迫辍学。之后，他自学了一年，并未再接受任何正式教育。\(^\text{[5]: 51 }\)

赫维赛德的姨夫是查尔斯·惠斯通爵士，这位国际知名的电报与电磁学专家是19世纪30年代中期首个商业化电报系统的共同发明者。惠斯通非常关心外甥的教育，\(^\text{[6]}\)并于1867年将赫维赛德送往英国纽卡斯尔（，在那里，他在惠斯通的哥哥亚瑟·惠斯通管理的一家电报公司任职。\(^\text{[5]: 53 }\)

两年后，赫维赛德在丹麦“大北电报公司”找到了一份电报操作员的工作，负责从纽卡斯尔铺设通往丹麦的海底电缆，施工方为英国承包商。不久之后，他被提升为电气工程师。在工作期间，赫维赛德仍坚持自学。22岁时，他在著名的《哲学杂志》上发表了论文《惠斯通电桥在给定电流计与电池条件下测量特定电阻的最佳配置》\(^\text{[7]}\)。这篇论文得到了多位物理学家的高度评价——其中包括威廉·汤姆孙爵士（Sir William Thomson，后来的开尔文勋爵）和詹姆斯·克拉克·麦克斯韦，他们此前都曾尝试解决这一代数问题但未成功。赫维赛德还亲自将论文送给了汤姆孙爵士。此后，他又发表了一篇关于“双工传输”（在电报电缆中的应用的论文\(^\text{[8]}\)，并在文中讽刺了当时英国邮政电报系统总工程师R.S.卡利，后者一直否认双工电报的可行性。1873年晚些时候，赫维赛德申请加入“电报工程师学会”，却被拒绝，理由是“他们不希望接收电报职员”。这一侮辱激怒了赫维赛德，他于是请威廉·汤姆孙为其作推荐人，并在该学会会长的支持下，最终“尽管邮局里那些势利小人反对”，他还是成功入会。\(^\text{[5]: 60 }\)

1873年，赫维赛德读到了刚刚出版、后来享誉世界的詹姆斯·克拉克·麦克斯韦两卷本巨著《电与磁论》。多年以后，他在晚年回忆道：

“我还记得年轻时第一次看到麦克斯韦那部伟大的著作……我立刻意识到它是伟大的、更加伟大的、最伟大的，其力量中蕴藏着惊人的潜能……我下定决心要精通这本书，并立即开始钻研。当时我非常无知，对数学分析一无所知（只学过一些中学代数和三角学，还大多忘了），因此我知道自己该怎么补课。我花了好几年时间，才尽我所能去理解其中的内容。之后，我暂且放下麦克斯韦，按照自己的方式继续研究。那时我的进步要快得多……可以说，我所宣扬的，是我自己对麦克斯韦理论的理解与诠释。”\(^\text{[8]}\)

赫维赛德在家中独自从事研究，发展出了传输线理论（transmission line theory，也称为“电报方程”）。他通过数学推导证明：在电报线路中，如果沿线均匀分布电感，就可以减小信号的衰减和失真；并且，如果电感足够大且绝缘电阻不过高，那么电路将不会发生频率失真——也就是说，各频率的电流将以相同的速度传播。\(^\text{[10]}\)赫维赛德的方程为电报系统的改进和推广提供了重要的理论基础。
\subsection{中年时期}
从1882年至1902年（中间有三年例外），赫维赛德定期为行业刊物《电气工程师》撰写文章。该杂志希望借助他的学术声望提升自身影响力，因此每年支付他40英镑稿酬。虽然这笔钱几乎不足以维生，但赫维赛德生活要求极低，而且这份工作正是他最热爱的事情。1883年至1887年间，他平均每月发表2至3篇论文，这些文章后来构成了他的重要著作《电磁理论》和《电学论文集》的主要内容。\(^\text{[5]: 71 }\)

1880年，赫维赛德研究了电报传输线中的“集肤效应”，同年他在英国申请了“同轴电缆”的专利。1884年，他将麦克斯韦原本复杂的数学表达式（当时曾以四元数形式改写）重新整理为现代的“向量表示法”，将原先20个未知数的20个方程大大简化为如今广为人知的4个偏微分方程，也就是我们今天所说的“麦克斯韦方程组”。这四个方程系统地描述了电荷（静态与动态）、磁场及其相互关系——即电磁场的本质。

1880年至1887年间，赫维赛德发展了“算子演算”，以符号$p$（布尔此前使用$D$表示\(^\text{[11]}\)）代表微分算子，从而可以将微分方程直接转化为代数方程求解。这种方法后来因“缺乏数学严谨性”而引发广泛争议。赫维赛德对此曾著名地辩解道：“数学是一门实验科学，定义并不是先行的，而是在研究对象的本质逐渐显现后才自然形成的。”\(^\text{[12]}\)他还幽默地反问过一句：“难道仅仅因为我不能完全理解消化的过程，就要拒绝吃饭吗？”\(^\text{[13]}\)1887年，赫维赛德与他的哥哥亚瑟合作撰写了一篇题为《桥式电话系统》的论文。然而，这篇论文被亚瑟的上司——英国邮政局的威廉·亨利·普里斯阻止发表。原因是论文中提出应在电话和电报线路中加入“负载线圈”（loading coils，即电感器），以增强自感，从而校正信号传输中的失真。而普里斯此前刚刚宣称“自感是清晰传输的最大敌人”，因此他坚决反对。\(^\text{[14]}\)赫维赛德还确信，正是普里斯在背后操纵，迫使《电气工程师》的编辑被解雇，导致他连载多年的论文系列被迫中断（直到1891年才恢复）。赫维赛德与普里斯之间的敌意由来已久。赫维赛德认为普里斯在数学上“完全不懂行”，传记作者保罗·J·纳欣也持相同观点：“普里斯是一个掌握权力的政府官员，野心勃勃，但在某些令人惊讶的方面却是个十足的笨蛋。” 普里斯阻挠赫维赛德工作的动机，更多是为了维护自己的声誉、避免承认错误，而非出于对赫维赛德研究本身的质疑。\(^\text{[4]: xi–xvii, 162–183 }\)

赫维赛德的研究成果在发表后很长一段时间内都未被重视。直到1897年，美国电话电报公司才开始关注这项工作，并聘请公司科学家乔治·A·坎贝尔以及外部研究员迈克尔·I·普平对赫维赛德的理论进行验证，试图找出其中的不足或错误。坎贝尔和普平在赫维赛德的基础上进行了扩展研究，AT&T随即申请了多项专利，不仅包括他们自己的研究成果，还涵盖了赫维赛德此前发明的线圈构造方法。后来，AT&T向赫维赛德提出金钱补偿，以换取其相关权利。据说这一提议部分源于贝尔系统工程师们对赫维赛德的尊敬。然而，赫维赛德拒绝了这一提议，表示除非公司公开承认其发明，否则他不会接受任何报酬。尽管他长期生活贫困，但他仍坚定地拒绝了这笔钱——这一举动尤为令人钦佩。1959年，诺伯特·维纳发表了小说《诱惑者》，在书中影射性地指控AT&T（化名为“威廉斯控制公司”Williams Controls Company）和迈克尔·I·普平（化名为“迭戈·多明格斯”Diego Dominguez）窃取了赫维赛德的发明成果。\(^\text{[15][16][17]}\)

然而，这次挫折使赫维赛德将注意力转向了电磁辐射。\(^\text{[18]}\)在他于1888年和1889年发表的两篇论文中，他计算了运动电荷周围电场与磁场的变形，以及电荷进入更高密度介质时所产生的效应。这项研究预言了后来被称为契伦科夫辐射的现象，并启发了他的好友乔治·菲茨杰拉德提出了如今所知的洛伦兹–菲茨杰拉德收缩假设。

1889年，赫维赛德首次发表了运动带电粒子所受磁力的正确推导，\(^\text{[19]}\)这正是现在所称洛伦兹力的磁场分量部分。

在19世纪80年代末至90年代初，赫维赛德进一步研究了电磁质量的概念。他认为电磁质量可被视为一种“物质质量”，能够产生相同的物理效应。后来，威廉·维恩在低速条件下验证了赫维赛德的公式。

1891年，英国皇家学会以授予他院士称号的方式，正式表彰他在电磁现象数学化描述方面的杰出贡献。次年，《皇家学会哲学汇刊》专门刊登了超过五十页篇幅的论文，介绍他的向量分析方法与电磁理论。

1894年，他被选为曼彻斯特文学与哲学学会的名誉会员。\(^\text{[20]}\)1905年，德国哥廷根大学授予他荣誉博士学位，以表彰其在电磁学方面的卓越成就。
\subsection{晚年与思想观点}

1896年，乔治·菲茨杰拉德和约翰·佩里为赫维赛德争取到了一项政府养老金，每年金额为120英镑。当时他已定居在德文郡。他们劝说他接受这笔资助，因为赫维赛德此前曾拒绝过皇家学会提供的其他慈善援助。\(^\text{[18]}\)

1902年，赫维赛德提出了如今被称为肯内利–赫维赛德层的电离层概念。他的理论解释了无线电信号能够绕地球曲率传播的原理。电离层的存在后来于1923年得到证实。赫维赛德的预言，加上普朗克辐射理论，可能使人们在此后的几十年中，对探测来自太阳及其他天体的无线电波失去了兴趣。无论出于何种原因，之后的三十年间几乎无人尝试，直到1932年詹斯基开创了射电天文学。

赫维赛德是阿尔伯特·爱因斯坦相对论的公开反对者。\(^\text{[21]}\)数学家霍华德·伊夫斯评论说，赫维赛德“是当时唯一一个质疑爱因斯坦的顶尖物理学家，他对相对论的批评常常近乎荒谬”。\(^\text{[21]}\)
